\section{How this book is structured}
\label{sec:basics:structure}

This book is two books, and the colour of the page says which one is in hand.

The first part, on white pages, is a course. Its chapters each gather a set of
related notions, and each notion appears at the point where it is first needed,
so the chapters are meant to be read in order. Every one of them ends with
exercises and their solutions.

The second part, on light grey pages, is the specification. It runs through the
same material in the same order, chapter for chapter, but formally and
exhaustively. It is for consulting rather than reading: a question about a notion
is answered in the part II chapter that matches the part I chapter where the
notion was taught.

Source code appears in listings. A badge in the top right corner of a listing
names the language it is written in, and the colour of its background says the
same thing again.

\begin{description}
\item[\ymirbadgeymir] \fcolorbox{black!35}{background!40}{\phantom{MM}}~a
  \textbf{yellow} background, for source code written in the Ymir language. This
  is by far the most common listing of the book.
\item[\ymirbadgemyil] \fcolorbox{black!35}{teal!10}{\phantom{MM}}~a
  \textbf{teal} background, for M-YIL code, and
  \fcolorbox{black!35}{olive!10}{\phantom{MM}}~a \textbf{green} one, for L-YIL
  code. These two languages are not written by the programmer, they are
  generated by the compiler, and are used in this book to describe what a Ymir
  program actually does once compiled. They are presented in
  Section~\ref{sec:basics:yil_language}.
\item[\ymirbadgeshell] \fcolorbox{black!35}{gray!10}{\phantom{MM}}~a
  \textbf{gray} background, for a terminal session. The lines starting with a
  prompt of the form \tokennolst{alice@dev:\mtildeascii{}\$} are commands typed
  by the user, and the lines that follow them are the output of those commands.
\item[\ymirbadgec] source code written in the \textbf{C} language, on the same
  background as an Ymir listing. It appears only where a program has to be
  linked against code that is not written in Ymir.
\end{description}

Some listings show a program that is wrong on purpose, so that the reason the
compiler rejects it can be explained. Such a listing is Ymir like any other, but
its badge is a warning sign, \ymirbadgeerror. Two highlights then appear inside
it: \hb{red} marks the code at fault, and \hcb{green} marks a correction applied
to it.

\vfill%
\pagebreak%
